


\section{Fuentes de obtención de enzimas}



\section{Localización celular de las enzimas}



\section{Métodos de extracción de enzimas}



\section{Métodos de caracterización de enzimas}



\section{Variables para medir actividad enzimática}

Existen diferentes variables que se pueden medir de una reacción para de terminar su velocidad:

\begin{itemize}
    \item Sustrato \cite{malairuangHighCellDensity2020}
    \item Producto
        \begin{itemize}
            \item Por ejemplo, aquí se podría medir la absorbancia del producto ($O$), cuyo valor, al cambiar en el tiempo, se puede usar para determinar la cinética de la reacción.
                \begin{align*}
                H_2O_2 + OH_2 &\Longrightarrow 2H_20 + O
                \\
                O \lambda &= 440nm
                \end{align*}
        \end{itemize}
    \item Cofactor
        \begin{itemize}
            \item En este ejemplo, tanto el $NAD^+$ como el $NADH$ funcionan como cofactores, y se pueden usar como elementos a medir para determinar la cinética de la reacción.
                \begin{align*}
                SH_2 + NAD^{+} &\Longrightarrow NADH + H^{+} + S
                \\
                NADH \lambda &= 340nm
                \\
                \\
                SH_2 + NADH &\Longrightarrow NAD^{+} + SH_2
                \\
                NAD^{+} \lambda &= 254NM
                \end{align*}
                
        \end{itemize}
    \item Reacción enzimática acoplada
    \item Reacción química acoplada
\end{itemize}


\section{Métodos de medición de actividad enzimática}

Existen diversas 2 formas de medir la actividad de una enzima: de manera directa y de manera indirecta.\cite{somboonSimpleGreenMethod2022}

La primera se utiliza cuando es posible realizar una medición de alguno de los siguientes aspectos:

\begin{itemize}
    \item La disminución del sustrato conforme avanza la reacción.
    \item La acumulación del producto conforme avanza la reacción.
\end{itemize}

Por otro lado, el método indirecto se realiza cuando no puede existir una medición directa de alguno de los siguientes aspectos:

\begin{itemize}
    \item Ni sustrato
    \item Ni complejo enzima-sustrato
    \item Ni producto
\end{itemize}

Para realizar este tipo de medición se pueden realizar las siguientes operaciones:

\begin{itemize}
    \item Medir el consumo de cofactor
    
    \item Usar una reacción acoplada que si se pueda medir. Esta puede ser una reacción enzimática o química y necesita de especificidad y muestras de control.
\end{itemize}

Las mediciones también se pueden clasificar en base a qué medida de tiempo se utilize:

\begin{itemize}
    \item A tiempo final: En esta se establece un tiempo determinado y se mide unicamente en $t_i$ y $t_f$. Por ejemplo, dada una reacción con durción de 10 minutos, solo se mediría en el minuto 0 (al inicio) y al minuto 10 (al final).\cite{lopezComparacionMetodosQue2017}
    \item Cinética: En esta se establece una unidad de tiempo que determina la periodicidad con que se medirá la reacción. Por ejemplo, medir cada 30 segundos. 
\end{itemize}

\section{Espectrofotometría en la medición de actividad enzimática}

\begin{equation}
    A = \varepsilon b C
\end{equation}

Donde:

\begin{itemize}
    \item $A$ = Absorbancia (Luz absorbida)
    \item $\varepsilon$ = Absortividad (Capacidad de la sustancia para absorber luz) 
    \item $b$ = Longitud de la celda
    \item $C$ = Concentración de la solución
\end{itemize}

\section{Ejemplo - Producción de melanina}\footnote{PENDIENTE DE MEJORA}

\begin{align*}
    \text{Tirosina} + \text{Tirosinasa (PPO)} + O_2 + Cu^{2+} &\Longrightarrow \text{Tirosina Oxidada} + DOPA + H_2O
    \\
    \\
    DOPA + PPO + O_2 &\Longrightarrow \text{DOPA-Quinona}
    \\
    \\
    &\Longrightarrow \text{[Reacciones no enzimáticas]}
    \\
    \\
    \text{DOPA-Quinona} &\Longrightarrow \text{Dopacromo (Naranja)}
    \\
    \\
    \text{Dopacromo} &\Longrightarrow \text{5-6-hidroxi-indol (Incoloro)} + H_2O
    \\
    \\
    \text{5-6-hidroxi-indol} &\Longrightarrow \text{Indol-5-6-Quinona (Amarillo)}
    \\
    \\
    \text{Indol-5-6-Quinona} &\Longrightarrow \text{Melanocromo (Púrpura)}
    \\
    \\
    \text{N-Melanocromo} &\Longrightarrow \text{Melanina (Negra)}
\end{align*}

En este ejempo, se puede medir el cambio de color en una sustancia que contenga melanina (frutas o verduras, por ejemplo), y a partir de ese cambio de color determinar, de manera general y tal vez no tan exacta, si existe actividad enzimática, pues aunque las reacciones que muestran cambio de color no son enzimáticas dependen de un producto que se genera solamente por actividad de enzimas.


% --------------------------------------------------------------------------------------------------- %

\newpage \putbib